%----------------------------------------------------------------------
\section{Directed Higher Categories}\label{sec:directed-higher}
%----------------------------------------------------------------------

Section \ref{sec:homotopy} constructed a symmetric $n$-groupoid
from the exterior algebra. Section \ref{sec:cattheory} constructed
directed 1-categories by layering $\tau/\sigma$ asymmetry on top.
We now show how these combine: the framework produces
$(n,1)$-categories, and obtaining higher directed structure
requires explicit meta-discrimination.

\subsection{Why direction is confined to grade 1}

\begin{theorem}[Natural $(n,1)$-structure]\label{thm:n1cat}
The category $\mathcal{C}(S,\kappa)$ with the
discrimination complex $\bigwedge(\GF(2)^n)$ is an
$(n,1)$-category: 1-morphisms are directed (potentially
non-invertible), and $k$-morphisms for $k \geq 2$ are
symmetric and invertible.
\end{theorem}

\begin{proof}
The argument turns on where $\tau/\sigma$ labeling applies.

\emph{Grade 1: directed.} A discriminator $d \in \bigwedge^1 V$
acts directly on the population $S$, producing assignments
$\tau_d(a), \sigma_d(a) \in \GF(2)$ for each $a \in S$.
The labels $\tau$ and $\sigma$ are structurally distinct
(Definition \ref{def:discriminator}), so $d$ induces a
directed morphism $a \to b$ when $\tau_d(a) = 1$ and
$\sigma_d(b) = 1$ (Definition \ref{def:directed-morphism}).
The reverse morphism requires a separate discriminator
(Proposition \ref{prop:noninvert}). Therefore 1-morphisms
are not automatically invertible.

\emph{Grade 2 and above: symmetric.}
A $k$-cell $\omega = d_{i_1} \wedge \cdots \wedge d_{i_k}
\in \bigwedge^k V$ for $k \geq 2$ does not act directly
on the population. It relates $k$ discriminators. The
$\tau/\sigma$ labeling is a property of the interface
between a discriminator and a population; since $\omega$
does not interface with $S$ directly, no $\tau/\sigma$
assignment applies to $\omega$. The $k$-cell inherits
only the algebraic properties of $\bigwedge(\GF(2)^n)$:
\begin{itemize}[nosep]
  \item $\omega \wedge \omega = 0$ (self-inverse),
  \item $\omega_1 \wedge \omega_2 = \omega_2 \wedge \omega_1$
    (symmetric),
\end{itemize}
so every $k$-morphism for $k \geq 2$ is invertible (its
own inverse) and composition is symmetric. These are exactly
the conditions for the $k$-morphisms to form a groupoid at
each level $k \geq 2$.

Combining: directed 1-morphisms with symmetric, invertible
$k$-morphisms for $k \geq 2$ is the definition of an
$(n,1)$-category (truncated at grade $n$).
\end{proof}

\begin{remark}
The $(n,1)$-structure is not a design choice. It follows from
the architecture: $\tau/\sigma$ labels the population;
grade 1 elements contact the population; grade $\geq 2$
elements do not. Direction exists exactly where the algebra
meets the world.
\end{remark}

\subsection{2-morphisms as homotopies between directed paths}

\begin{definition}[Directed homotopy]\label{def:directed-homotopy}
Let $f, g: A \to B$ be two directed morphisms witnessed by
discriminators $d_f$ and $d_g$ respectively. A
\emph{2-morphism} $\alpha: f \Rightarrow g$ is the 2-cell
$d_f \wedge d_g \in \bigwedge^2 V$. This 2-cell is:
\begin{itemize}[nosep]
  \item symmetric: $d_f \wedge d_g = d_g \wedge d_f$, so
    $\alpha: f \Rightarrow g$ and $\alpha: g \Rightarrow f$
    are the same 2-cell;
  \item self-inverse: $(d_f \wedge d_g) \wedge (d_f \wedge d_g) = 0$;
  \item well-bounded: $\partial(d_f \wedge d_g) = \{d_f, d_g\}$
    (Definition \ref{def:boundary}).
\end{itemize}
\end{definition}

\begin{proposition}[2-morphisms witness
  indiscernibility of paths]\label{prop:2morph}
The 2-cell $d_f \wedge d_g$ exists (is non-zero) if and
only if $d_f \neq d_g$ (the two discriminators are
linearly independent). When it exists, it witnesses that
$f$ and $g$ are two \emph{distinct} directed paths from
$A$ to $B$, and provides a symmetric, invertible
identification between them.

When $d_f = d_g$, the 2-cell is $d_f \wedge d_f = 0$:
the degenerate case. Identical morphisms are connected
by the trivial 2-morphism. This is the identity 2-cell.
\end{proposition}

\begin{proof}
Linear independence of $d_f$ and $d_g$ is equivalent to
$d_f \wedge d_g \neq 0$ in $\bigwedge^2(\GF(2)^n)$.
If $d_f = d_g$, then $d_f \wedge d_g = d_f \wedge d_f = 0$.
The boundary $\partial(d_f \wedge d_g) = \{d_f, d_g\}$
confirms that the 2-cell connects exactly the two
1-cells that constitute its factors.
\end{proof}

\begin{remark}
The symmetry of the 2-cell means there is no preferred
direction of the homotopy between $f$ and $g$. The
two directed paths are identified without one being
``before'' or ``after'' the other. This is the correct
behavior for an $(n,1)$-category: 2-morphisms are
equivalences, not directed transformations.
\end{remark}

\begin{example}[Directed 2-morphism: two classification schemes]
\label{ex:2morph}
Consider a population $S$ of animals. Two discriminators
$d_1$ and $d_2$ both classify animals as ``dangerous''
($\tau$) or ``safe'' ($\sigma$):
\begin{itemize}[nosep]
  \item $d_1$: ``has venom'' (a biochemical discriminator).
    Snake $\to \tau$, rabbit $\to \sigma$.
  \item $d_2$: ``is a predator'' (an ecological discriminator).
    Snake $\to \tau$, rabbit $\to \sigma$.
\end{itemize}

Both $d_1$ and $d_2$ are directed 1-morphisms: they send
``snake'' from $\sigma$ to $\tau$ (classifying it as
dangerous). They agree on the snake and the rabbit. But
they are distinct discriminators (linearly independent in
$\kappa$), so the 2-cell $d_1 \wedge d_2 \neq 0$.

This 2-cell is the directed 2-morphism witnessing that
``venomous'' and ``predatory'' are two \emph{distinct}
paths from ``unclassified animal'' to ``dangerous.'' The
2-cell does not prefer one path over the other (it is
symmetric: $d_1 \wedge d_2 = d_2 \wedge d_1$). It
records that two independent criteria happen to agree on
this population --- and that they \emph{could} disagree
on a different population (e.g.\ a bear: predatory but
not venomous). The 2-cell's existence anticipates the
residue that a richer population would produce.

Under the $\sigma$-perspective, the 2-cell becomes a
\emph{directed} 2-morphism: a discrimination between the
two classification schemes themselves. The
$\sigma$-discriminator asks: ``does this element's
classification change when we switch from biochemical to
ecological criteria?'' For the snake: no. For the bear:
yes. The directed 2-morphism is the structural record of
this inter-scheme comparison.
\end{example}

\subsection{Higher cells and the homotopy tower}

\begin{proposition}[Structure of $k$-morphisms]\label{prop:k-morph}
For $k \geq 2$, a $k$-cell
$\omega = d_1 \wedge \cdots \wedge d_k$ witnesses the
joint relationship among $k$ directed 1-morphisms. Its
boundary $\partial \omega$ consists of $k$ distinct
$(k-1)$-cells (Proposition \ref{prop:boundary}), each
obtained by omitting one factor. The $k$-cell asserts
that these $(k-1)$-cells are coherently related.

The tower of $k$-morphisms for $k \geq 2$ forms a
symmetric groupoid at each level (Theorem \ref{thm:groupoid}).
The entire tower is the homotopy type of the space of
directed paths between two objects, computed over $\GF(2)$.
\end{proposition}

\begin{remark}[Truncation level as complexity measure]
In the $(n,1)$-category $\mathcal{C}(S,\kappa)$ with
$\dim(\kappa) = n$, the highest non-trivial $k$-morphism
is at grade $n$. This is the \emph{top cell}
$d_1 \wedge d_2 \wedge \cdots \wedge d_n \in \bigwedge^n V$,
the unique (up to scalar) element of maximum grade. It
witnesses the coherence of \emph{all} discriminators
simultaneously --- the total relationship among every
discrimination that has been performed. As $\kappa$ evolves,
the truncation level rises, and deeper coherence relations
become expressible.
\end{remark}

\subsection{The triangle identity and free meta-discrimination}
\label{sec:meta}

The $(n,1)$-structure of Theorem \ref{thm:n1cat} assumed
that $\tau/\sigma$ labeling exists only at the interface
between discriminators and population. We now show that
meta-discrimination --- labeling of discriminators
themselves --- is already present in the algebraic
structure, requiring no additional construction.

\begin{definition}[Triangle identity]\label{def:triangle}
The discrimination regime $\kappa$, together with its
membership masks $\sigma$ and $\tau$, satisfies the
\emph{triangle identity}:
\begin{equation}\label{eq:triangle}
  \kappa = \sigma \otimes \tau.
\end{equation}
Over $\GF(2)$, this identity has a triadic symmetry.
Since $\kappa = \sigma + \tau$ (in $\GF(2)$, where
$\otimes$ reduces to XOR), we also have:
\begin{align}
  \sigma &= \kappa + \tau, \label{eq:tri-sigma}\\
  \tau   &= \kappa + \sigma. \label{eq:tri-tau}
\end{align}
Each of the three elements is the XOR of the other two.
No element is privileged. Any of $\kappa$, $\sigma$, $\tau$
can serve as the discrimination regime, with the remaining
two as membership masks.
\end{definition}

\begin{remark}[Grounding the identification]
\label{rem:triangle-grounding}
The triangle identity identifies three objects that
appear to be different kinds of things: $\kappa$ is a
regime (an ordered sequence of discriminators with
monotonic growth, a write path, and a read path), while
$\sigma$ and $\tau$ are masks (functions from elements
to $\GF(2)$). The identification proceeds as follows.

Each discriminator $d \in \kappa$ produces a pair
$(\tau_d, \sigma_d)$ of functions $S \to \GF(2)$. The
XOR signature $\chi_d = \tau_d + \sigma_d$ is a single
function $S \to \GF(2)$. The regime $\kappa$ is thus
associated with a family of $\GF(2)$-valued functions
$\{\chi_{d_1}, \ldots, \chi_{d_n}\}$. Each $\chi_{d_i}$
is an element of $\GF(2)^S$ (the vector space of
functions $S \to \GF(2)$).

The masks $\sigma$ and $\tau$ are also elements of
$\GF(2)^S$ (via $\tau_d: S \to \GF(2)$ and
$\sigma_d: S \to \GF(2)$). The identity
$\chi_d = \tau_d + \sigma_d$ is an equation in
$\GF(2)^S$: it holds pointwise, for every element
of $S$.

The $S_3$ symmetry acts on this equation by permuting
which function plays which role: the XOR signature (which
records whether exactly one predicate fires), the
$\tau$-comprehension (which records whether $\tau$ fires),
and the $\sigma$-comprehension (which records whether
$\sigma$ fires). All three are elements of the
same vector space $\GF(2)^S$. The structural properties
attached to $\kappa$ --- monotonic growth, the write path,
the read path --- transfer under permutation because they
are consequences of the $\GF(2)$ algebra, not of the
label (Proposition \ref{prop:dynamics-transfer}).
\end{remark}

\begin{proposition}[Three perspectives]\label{prop:perspectives}
The triangle identity yields three equally valid, equally
lossless perspectives on the same discrimination structure:
\begin{enumerate}[label=(\roman*),nosep]
  \item \textbf{$\kappa$-perspective}: $\kappa$ is the
    regime, $\sigma$ and $\tau$ are masks.
    Discriminators act on the population. This is the
    standard framework (Section \ref{sec:framework}),
    giving directed 1-morphisms.
  \item \textbf{$\sigma$-perspective}: $\sigma$ is the
    regime, $\kappa$ and $\tau$ are masks.
    Discriminators of the $\sigma$-regime act on the
    discriminators of $\kappa$ (treated as a population).
    This gives directed 2-morphisms: directed
    relationships between 1-morphisms.
  \item \textbf{$\tau$-perspective}: $\tau$ is the regime,
    $\kappa$ and $\sigma$ are masks.
    This provides a third directed perspective,
    independently classifying at yet another level.
\end{enumerate}
Each perspective is an equally lossless view of the
underlying structure. No information is created or
destroyed by changing perspectives.
\end{proposition}

\begin{proof}
The triangle identity (\ref{eq:triangle})--(\ref{eq:tri-tau})
is symmetric under any permutation of $\{\kappa, \sigma, \tau\}$.
Given any two of the three, the third is determined (their
XOR). Therefore each perspective contains exactly the same
information as the others. The map between perspectives is
its own inverse (XOR is involutive), so no perspective is
lossy.
\end{proof}

\begin{proposition}[Dynamics transfer under permutation]
\label{prop:dynamics-transfer}
The structural properties of the $\kappa$-perspective ---
monotonic growth (Proposition \ref{prop:monotone}), the
linear write path (Definition \ref{def:evolution}), and
the monoidal read path (Proposition \ref{prop:monoidal})
--- transfer to the $\sigma$- and $\tau$-perspectives
under $S_3$ permutation.
\end{proposition}

\begin{proof}
These properties are consequences of the $\GF(2)$
algebra, not of the label ``$\kappa$.''

\emph{Monotonicity.} $\kappa$-evolution is monotonic
because each evolution adds a linearly independent basis
vector to $V = \GF(2)^n$, increasing $n$ by one
(Proposition \ref{prop:monotone}). Under the
$\sigma$-perspective, $\sigma$ is the regime and its
``discriminators'' are the $\sigma$-masks viewed as a
population. Evolution in the $\sigma$-perspective adds a
new element to this population --- equivalently, a new
dimension in $\GF(2)^n$ (since $\sigma = \kappa + \tau$,
a new $\kappa$-dimension produces a new $\sigma$-dimension
by the triangle identity). The growth is monotonic for the
same algebraic reason: dimensions of a vector space over
$\GF(2)$ grow monotonically under basis extension.

\emph{Write/read.} The write path's linearity (each
residue consumed once) is a property of the exterior
algebra's nilpotency: $d \wedge d = 0$ prevents a
discriminator from being consumed twice. This holds for
any element of $\bigwedge V$, regardless of whether it is
labeled $\kappa$, $\sigma$, or $\tau$. The read path's
monoidal property (order-independence of filtration) is
Proposition \ref{prop:monoidal}, which depends on the
independence of discriminators' profiles --- again an
algebraic property that holds regardless of labeling.

Therefore: when $\sigma$ plays the role of regime, it
inherits monotonic growth, a linear write path, and a
monoidal read path. The same holds for $\tau$. The
dynamics are properties of the algebra, not of the label.
\end{proof}

\begin{remark}[Semantic asymmetry across perspectives]
Proposition \ref{prop:perspectives} proves that algebraic
properties (monotonicity, nilpotency, order-independence)
transfer under $S_3$ rotation. But when $\sigma$ plays
the role of regime, the ``population'' being discriminated
is the set of discriminators of $\kappa$, which carry
rich internal structure (the vector space structure of $V$,
profile linearity (Definition \ref{def:profile-linear}), evaluation linearity (Definition \ref{def:eval-linear})). The rotation is
algebraically valid, but the semantic content at different
levels may differ: discriminators-as-data are not generic
data. The paper treats the rotation as formal; a fuller
treatment would verify that the additional structure of
discriminators-as-data does not invalidate the algebraic
transfer.
\end{remark}

\begin{definition}[$\GF(2)$ toroid]\label{def:toroid}
The three elements $\kappa, \sigma, \tau$ are the three
non-zero points of $\GF(2)^2$:
\[
  \kappa \leftrightarrow (1,1), \qquad
  \sigma \leftrightarrow (0,1), \qquad
  \tau   \leftrightarrow (1,0).
\]
(Any assignment satisfying the pairwise-XOR condition is
valid.) The space $\GF(2)^2$ is a \emph{discrete torus}
$\mathbb{Z}/2\mathbb{Z} \times \mathbb{Z}/2\mathbb{Z}$,
and the three non-zero points form a triangle within it.
The symmetry group of this triangle is $S_3$, the
symmetric group on three elements, acting by permutation
of $\{\kappa, \sigma, \tau\}$.

Note: $\GF(2)^2 = \GF(2) \times \GF(2)$ is the Klein
four-group. The product structure is the paper's
central operation: composing two $\GF(2)$-valued
discriminator evaluations into a $\GF(2)^2$-valued
membership profile. Conversely, $\GF(2)$ is the
quotient $\GF(2)^2 / \GF(2)$: the XOR signature
$\chi = \tau + \sigma$ is the quotient map, and the
residue $R = \{a : \chi(a) = 0\}$ is the kernel.
The framework systematically traverses this
product-and-quotient structure: upward from $\GF(2)$
to $\GF(2)^2$ when composing discriminator pairs,
downward from $\GF(2)^2$ to $\GF(2)$ when computing
the XOR signature. The ``toroid'' names this
bidirectional traversal, not merely the group.
\end{definition}

\begin{theorem}[Free $(n,k)$-structure]\label{thm:nk-free}
The $S_3$ symmetry of the triangle identity provides
directed structure at multiple levels without constructing
meta-regimes. Each element of $S_3$ yields a perspective
in which a different component plays the role of
discrimination regime, providing direction at the
corresponding categorical level. Specifically:
\begin{itemize}[nosep]
  \item The three transpositions (swaps of two elements)
    each exchange which component is ``regime'' vs.\
    ``mask,'' providing directed structure at a
    different level.
  \item The two 3-cycles rotate all three roles, providing
    composite perspectives that chain through all levels.
  \item The identity preserves the standard $\kappa$-perspective.
\end{itemize}
The resulting structure is an $(n,k)$-category where the
number of directed levels $k$ is determined by the
subgroup of $S_3$ that is actively employed.
\end{theorem}

\begin{proof}
Each perspective (Proposition \ref{prop:perspectives})
provides $\tau/\sigma$ labeling at one categorical level:
the level at which the ``population'' of that perspective
lives. The $\kappa$-perspective labels elements of $S$
(level 0 entities), giving directed 1-morphisms. The
$\sigma$-perspective labels discriminators of $\kappa$
(level 1 entities), giving directed 2-morphisms. The
$\tau$-perspective labels discriminators of $\sigma$
(level 2 entities), giving directed 3-morphisms.

Since each perspective is equally lossless (Proposition
\ref{prop:perspectives}), the directed structure at each
level is as legitimate as at level 1. No additional
structure is constructed; the direction was always present
in the triadic symmetry and merely becomes visible upon
rotating the perspective.

For $k$ levels of direction, apply $k$ of the three
available perspectives. The maximum from the triangle
identity alone is $k = 3$ (all three perspectives active
simultaneously). For $k > 3$, iterate: apply the triangle
identity to the meta-regime itself, obtaining a tower of
nested triangles. Each iteration provides three additional
levels of direction.
\end{proof}

\begin{corollary}[Direction requires no additional axiom]\label{cor:free-direction}
Meta-discrimination does not cost additional structure.
It is a rotation of the existing $\GF(2)$ toroid. The
apparent ``cost'' of direction at higher levels was an
artefact of fixing the $\kappa$-perspective and treating
the other perspectives as external constructions. Under
the toroidal symmetry, all perspectives are internal.
\end{corollary}

\begin{remark}[Geometric content of the toroid]
The $\GF(2)$ toroid $\mathbb{Z}/2\mathbb{Z} \times
\mathbb{Z}/2\mathbb{Z}$ is the simplest non-trivial
compact surface over a finite field. The three non-zero
points form an equilateral structure (all pairwise
distances are equal under the Hamming metric). The
$S_3$ symmetry is the full symmetry group of this
configuration. Geometric manipulations of the toroid ---
rotations, reflections, translations --- correspond to
changes of categorical perspective.

This gives the discrimination framework a geometric
foundation: the space of all possible categorical
perspectives on a given discrimination structure is
the $\GF(2)$ toroid, and navigating between perspectives
is geometric manipulation of this space.
\end{remark}

\begin{remark}[Why $(\infty,1)$-categories appear natural]
The $(\infty,1)$-category remains the structure most
commonly encountered in mathematical practice, despite
higher directed structure being freely available. This
is because the $\kappa$-perspective --- discriminating
elements of a population --- is the perspective most
naturally adopted when reasoning about mathematical
objects. The $\sigma$- and $\tau$-perspectives
(discriminating discriminations) are available but less
frequently invoked. The prevalence of $(\infty,1)$-categories
reflects a conventional choice of perspective, not a
structural limitation.

When mathematical practice does invoke directed
2-morphisms (in 2-category theory, double categories,
bicategories, or enriched category theory), the framework
predicts that a rotation to the $\sigma$-perspective has
been implicitly performed. This rotation requires no
additional axioms and preserves algebraic properties
(Proposition \ref{prop:perspectives}), though the
semantic content at different levels may differ
(Remark above).
\end{remark}

\begin{remark}[The $(\infty,\infty)$-category]
By iterating the triangle identity (applying it to
meta-regimes, then meta-meta-regimes, and so on),
direction is available at every level without bound.
As $\dim(\kappa) \to \infty$ and the tower of
iterated triangle identities extends without bound,
the resulting structure is an $(\infty,\infty)$-category:
directed at every level, with no truncation on either
homotopical complexity or directed depth. By the
indistinguishability principle (Proposition
\ref{prop:infinity}), this is operationally identical
to a completed $(\infty,\infty)$-category --- the
most general categorical structure in the standard
hierarchy.
\end{remark}

%----------------------------------------------------------------------
\section{Monoidal Closed Structure and the Periodic Table}
\label{sec:monoidal}
%----------------------------------------------------------------------

We now show that the $n$-categorical structure of the
discrimination framework and monoidal structure are the
same thing viewed from different vertices of the $\GF(2)$
toroid. The rotation that turns $n$-categories into monoidal
$(n-1)$-categories is delooping, and it is a free geometric
operation on the toroid. Closure (internal hom) falls out of
the triangle identity.
